% ===========================================================================
%  Typography & layout for the PDF edition of
%  "Improving Your Statistical Inferences"
%
%  Engine: LuaLaTeX (required for OpenType Literata + its ligatures).
%  Set in _quarto.yml via  pdf-engine: lualatex.
%
%  Fonts are bundled in include/fonts/ so the book builds identically on any
%  machine or CI runner without a system font install. The only extra TeX
%  package needed beyond a default TinyTeX is the math font:
%      tlmgr install stix2-otf
% ===========================================================================

% ---- Fonts -----------------------------------------------------------------
\usepackage{fontspec}
\usepackage{unicode-math}

% Body: EB Garamond. An elegant Garamond revival (installed as a TeX package,
% so no bundled files needed). Old-style figures sit well in running prose.
% A slight upscale lifts Garamond's small x-height to a comfortable reading
% size next to the code and math.
\setmainfont{EB Garamond}[
  Numbers   = OldStyle ,
  Scale     = 1.05 ,
  Ligatures = { Common, Contextual }
]

% Math is set in STIX Two Math (EB Garamond has no bundled math companion;
% STIX Two pairs acceptably and covers the symbols used in the book).
\setmathfont{STIX Two Math}

% Monospace for code: Latin Modern Mono ships with every TeX install, so no
% extra font file is required.
\setmonofont{Latin Modern Mono}[Scale=0.86]

\usepackage{microtype}               % subtle protrusion polish (LuaLaTeX)
\usepackage{graphicx}                % \includegraphics (logo on the cover)
\usepackage{tikz}                    % full-bleed cover background / layout
\usepackage{float}                   % enables the [H] "here, exactly" float spec

% Figure backgrounds: plots set par(bg = backgroundcolor), and my_setup.R sets
% that to white (#ffffff) for LaTeX output (snow #fffafa is only for HTML). A
% full render (freeze removed) regenerates every figure on clean white, so no
% PDF-side backing trick is needed here.

% ---- Colours ---------------------------------------------------------------
\usepackage{xcolor}
\definecolor{brandred}{HTML}{A60A0A}   % heading / accent colour (site brand)
\definecolor{inkdark}{HTML}{1C1717}    % near-black body ink
\definecolor{linkcol}{HTML}{8A1B1B}    % restrained link colour (muted red)
\definecolor{rulecol}{HTML}{C9433B}    % chapter-number rule

% ---- Page geometry & spacing ----------------------------------------------
\usepackage{geometry}
\geometry{
  paper=a4paper,
  inner=3.0cm, outer=3.0cm,
  top=3.0cm, bottom=3.2cm,
  headsep=0.7\baselineskip
}
\usepackage{setspace}
\setstretch{1.06}                      % slightly open leading for readability

% ---- Running headers / footers --------------------------------------------
\usepackage{fancyhdr}
\usepackage{emptypage}                 % no headers on blank pages
\renewcommand{\headrulewidth}{0.4pt}
\renewcommand{\footrulewidth}{0pt}
\fancypagestyle{mainstyle}{%
  \fancyhf{}
  \fancyhead[LE]{\small\textsc{\nouppercase{\leftmark}}}
  \fancyhead[RO]{\small\textsc{\nouppercase{\rightmark}}}
  \fancyfoot[C]{\small\thepage}
  \renewcommand{\headrule}{\color{rulecol!55}\hrule height \headrulewidth}
}
\pagestyle{mainstyle}
% Plain pages (chapter openers) get just a centered page number.
\fancypagestyle{plain}{\fancyhf{}\fancyfoot[C]{\small\thepage}\renewcommand{\headrulewidth}{0pt}}

% ---- Chapter & section styling --------------------------------------------
\usepackage{titlesec}

% Headings are set in Open Sans (the brand's HTML heading font), bundled in
% include/fonts/ so the build stays reproducible on CI without a system-font
% install. The sans gives strong contrast against the EB Garamond body and a
% modern-textbook feel; ExtraBold weight makes chapter/section titles command
% the page. (Open Sans is under the SIL OFL, like Literata — see
% include/fonts/OpenSans-LICENSE.txt.)
\newfontfamily\headingfont{Open Sans}[
  Path        = ./include/fonts/ ,
  Extension   = .ttf ,
  UprightFont = OpenSans-ExtraBold ,
  ItalicFont  = OpenSans-ExtraBoldItalic ,
  BoldFont    = OpenSans-ExtraBold ,
  BoldItalicFont = OpenSans-ExtraBoldItalic
]

% Chapter opener: large brand-red number on its own line, a thin rule,
% then the chapter title beneath it.
\titleformat{\chapter}[display]
  {\headingfont}
  {\flushright\fontsize{64}{64}\selectfont\color{brandred!85}\thechapter}
  {8pt}
  {\color{rulecol}\titlerule[1.2pt]\vspace{6pt}%
   \Huge\color{inkdark}\headingfont}
\titlespacing*{\chapter}{0pt}{8pt}{34pt}

\titleformat{\section}
  {\headingfont\Large\color{brandred!90}}{\thesection}{0.7em}{}
\titleformat{\subsection}
  {\headingfont\large\color{inkdark}}{\thesubsection}{0.6em}{}
\titleformat{\subsubsection}
  {\headingfont\normalsize\color{inkdark}}{\thesubsubsection}{0.5em}{}
% Once titlesec manages headings, every level used in the document needs a
% format or it errors ("No format for this command"). The changelog uses
% \paragraph, so define run-in formats for \paragraph and \subparagraph too.
\titleformat{\paragraph}[runin]
  {\headingfont\normalsize\color{inkdark}}{\theparagraph}{0.5em}{}
\titleformat{\subparagraph}[runin]
  {\headingfont\normalsize\color{inkdark}}{\thesubparagraph}{0.5em}{}
\titlespacing*{\section}{0pt}{1.5\baselineskip}{0.5\baselineskip}

% ---- Links -----------------------------------------------------------------
% Tame Quarto's default bright blue. (colorlinks is set in _quarto.yml.)
% Quarto loads hyperref later in the template, but include-in-header content is
% injected before that, so \hypersetup is not yet defined here. Defer our link
% colour settings until hyperref has been loaded.
\AtBeginDocument{\hypersetup{
  linkcolor=linkcol,
  citecolor=linkcol,
  urlcolor=linkcol,
  filecolor=linkcol
}}

% ---- Callout boxes (kept from the original, restyled to match) -------------
\usepackage{tcolorbox}
\tcbuselibrary{skins,breakable}

\definecolor{dangcol}{RGB}{152,62,130}
\definecolor{warncol}{RGB}{200,150,20}
\definecolor{infocol}{RGB}{70,122,172}
\definecolor{trycol}{RGB}{97,88,156}

\tcbset{
  calloutstyle/.style={
    enhanced, breakable, boxsep=5pt, arc=3pt, boxrule=0.9pt,
    left=8pt, right=8pt, top=6pt, bottom=6pt,
    coltext=black, fonttitle=\sffamily\bfseries
  }
}
\newtcolorbox{dangerous}{calloutstyle, colback=dangcol!8,  colframe=dangcol}
\newtcolorbox{warning}{calloutstyle,   colback=warncol!10, colframe=warncol!85!black}
\newtcolorbox{info}{calloutstyle,      colback=infocol!8,  colframe=infocol}
\newtcolorbox{try}{calloutstyle,       colback=trycol!8,   colframe=trycol}

% ---- Code blocks -----------------------------------------------------------
% Quarto highlights code in an fvextra Verbatim ("Highlighting") environment.
% By default long lines run off the right margin in the PDF. Turn on line
% breaking so code wraps within the text block, with a small arrow marking
% continued lines and breaks allowed at most non-word characters.
\usepackage{fvextra}
\fvset{
  breaklines=true,
  breakanywhere=true,
  breaksymbolleft=\raisebox{0.8ex}{\small\color{rulecol}$\hookrightarrow$\space},
  fontsize=\small
}

% Printed R *output* (results, console text) is emitted by Quarto in a plain
% {verbatim} environment, which does NOT wrap and so runs past the right
% margin (e.g. long "alternative hypothesis:" lines from t.test). Re-implement
% verbatim on top of fvextra's Verbatim so it inherits the breaklines settings
% above and wraps within the text block like the code does.
\AtBeginDocument{%
  \RecustomVerbatimEnvironment{verbatim}{Verbatim}{breaklines=true,breakanywhere=true,fontsize=\small}%
}

% ---- Block quotes ----------------------------------------------------------
% The HTML edition gives quotes a distinct, set-apart treatment; the book
% should be at least as explicit. Pandoc renders markdown ">" quotes as the
% {quote} environment, so we redefine it as a tinted tcolorbox with a thick
% brand-red bar down the left edge and slightly larger, italic text. This makes
% the many epigraph-style quotations in the book read clearly as quoted voice.
\definecolor{quotebg}{HTML}{FBF3F2}    % very pale warm tint behind quotes
\newtcolorbox{quotebox}{%
  enhanced, breakable, frame hidden, boxrule=0pt, arc=0pt, outer arc=0pt,
  colback=quotebg,
  borderline west={3pt}{0pt}{brandred},
  left=14pt, right=12pt, top=8pt, bottom=8pt,
  before skip=10pt, after skip=10pt
}
\renewenvironment{quote}
  {\begin{quotebox}\itshape\large\color{inkdark}\setstretch{1.12}}
  {\end{quotebox}}

% ---- Tables ----------------------------------------------------------------
\usepackage{booktabs}
\usepackage{colortbl}                  % \rowcolor / \arrayrulecolor
\renewcommand{\arraystretch}{1.12}     % a little more row breathing room

% Brand the kable/booktabs tables: a pale-red shaded header row, white bold
% header text on a stronger red is too heavy for data tables, so we use a soft
% tint with dark ink and recolour booktabs' rules in the brand red. Quarto's
% kable tables use \toprule/\midrule/\bottomrule, so colouring those rules
% restyles every table consistently.
\definecolor{tableheadbg}{HTML}{F2D9D7} % soft red tint for the header row
\arrayrulecolor{rulecol}
\setlength{\arrayrulewidth}{0.7pt}
% Shade the first row of every longtable/tabular header. Quarto emits the
% header cells right after \toprule; we hook \toprule to start the row colour
% and \midrule to end it, so the head band appears without editing each table.
\let\oldtoprule\toprule
\let\oldmidrule\midrule
\renewcommand{\toprule}{\oldtoprule\rowcolor{tableheadbg}}
\renewcommand{\midrule}{\rowcolor{white}\oldmidrule}

% ---- Table of contents tweak ----------------------------------------------
% Colour the chapter entries (KOMA-friendly; avoids tocloft which clashes).
\addtokomafont{chapterentry}{\color{inkdark}}

% ---- Title page -----------------------------------------------------------
% The cover (include/cover.tex, injected via include-before-body) is the front
% page. Quarto would otherwise emit its own plain \maketitle title page after
% it; suppress that so the cover is page 1 and is not duplicated.
\renewcommand{\maketitle}{}
